% =====================================================================
%  preamble.tex — Configuración global del documento
%  Trabajo Terminal 2026-B136 — ESCOM / IPN
%  Motor recomendado: LuaLaTeX  (lualatex / latexmk -lualatex)
% =====================================================================

% ---------- Idioma y codificación -----------------------------------
\usepackage[spanish,es-tabla,es-noindentfirst,es-nosectiondot]{babel}
\usepackage{csquotes}

% ---------- Tipografía (fontspec, LuaLaTeX) -------------------------
% Times New Roman en TODO el documento (cuerpo, títulos, portada,
% encabezados), igual que el documento Word original. Solo el código
% conserva una fuente monoespaciada (legibilidad/alineación).
\usepackage{fontspec}
\setmainfont{Times New Roman}      % cuerpo
\setsansfont{Times New Roman}      % títulos y encabezados (todo Times)
\setmonofont[Scale=0.88]{DejaVu Sans Mono} % solo bloques de código
\usepackage{microtype}             % refinamiento tipográfico

% ---------- Geometría y espaciado -----------------------------------
\usepackage[letterpaper,top=2.3cm,bottom=2.3cm,left=2.5cm,right=2.3cm]{geometry}
\usepackage{setspace}
\singlespacing                     % interlineado sencillo (igual que el Word)
\usepackage{ragged2e}
\usepackage[hyphens]{url}

% ---------- Matemáticas ---------------------------------------------
\usepackage{amsmath,amssymb,mathtools}

% ---------- Glifos Unicode ausentes en TeX Gyre (super/subíndices,
%            flechas) → se componen en modo matemático ---------------
\usepackage{newunicodechar}
\newunicodechar{⁰}{\ensuremath{^{0}}}\newunicodechar{¹}{\ensuremath{^{1}}}
\newunicodechar{²}{\ensuremath{^{2}}}\newunicodechar{³}{\ensuremath{^{3}}}
\newunicodechar{⁴}{\ensuremath{^{4}}}\newunicodechar{⁵}{\ensuremath{^{5}}}
\newunicodechar{⁶}{\ensuremath{^{6}}}\newunicodechar{⁷}{\ensuremath{^{7}}}
\newunicodechar{⁸}{\ensuremath{^{8}}}\newunicodechar{⁹}{\ensuremath{^{9}}}
\newunicodechar{⁺}{\ensuremath{^{+}}}\newunicodechar{⁻}{\ensuremath{^{-}}}
\newunicodechar{ⁿ}{\ensuremath{^{n}}}
\newunicodechar{₀}{\ensuremath{_{0}}}\newunicodechar{₁}{\ensuremath{_{1}}}
\newunicodechar{₂}{\ensuremath{_{2}}}\newunicodechar{₃}{\ensuremath{_{3}}}
\newunicodechar{₄}{\ensuremath{_{4}}}\newunicodechar{₅}{\ensuremath{_{5}}}
\newunicodechar{₆}{\ensuremath{_{6}}}\newunicodechar{₇}{\ensuremath{_{7}}}
\newunicodechar{₈}{\ensuremath{_{8}}}\newunicodechar{₉}{\ensuremath{_{9}}}
\newunicodechar{₊}{\ensuremath{_{+}}}\newunicodechar{₋}{\ensuremath{_{-}}}
\newunicodechar{↔}{\ensuremath{\leftrightarrow}}
\newunicodechar{→}{\ensuremath{\rightarrow}}\newunicodechar{←}{\ensuremath{\leftarrow}}
\newunicodechar{⇒}{\ensuremath{\Rightarrow}}\newunicodechar{⇔}{\ensuremath{\Leftrightarrow}}
\newunicodechar{≤}{\ensuremath{\leq}}\newunicodechar{≥}{\ensuremath{\geq}}
\newunicodechar{≈}{\ensuremath{\approx}}\newunicodechar{≠}{\ensuremath{\neq}}
\newunicodechar{×}{\ensuremath{\times}}\newunicodechar{∈}{\ensuremath{\in}}
\newunicodechar{‖}{\ensuremath{\Vert}}\newunicodechar{∥}{\ensuremath{\Vert}}

% ---------- Gráficos y figuras --------------------------------------
\usepackage{graphicx}
\graphicspath{{figures/}}
\usepackage{float}
\usepackage{subcaption}
\usepackage[font=small,labelfont=bf,margin=1cm]{caption}

% ---------- Tablas ---------------------------------------------------
\usepackage{booktabs}
\usepackage{longtable}
\usepackage{tabularx}
\usepackage{array}
\usepackage{multirow}
\usepackage{makecell}
\renewcommand{\arraystretch}{1.15}
% Columnas tipo X con alineaciones
\newcolumntype{L}[1]{>{\raggedright\arraybackslash}p{#1}}
\newcolumntype{C}[1]{>{\centering\arraybackslash}p{#1}}
\newcolumntype{R}[1]{>{\raggedleft\arraybackslash}p{#1}}

% ---------- Listas ---------------------------------------------------
\usepackage{enumitem}
\setlist{itemsep=2pt,topsep=4pt}

% ---------- Código ---------------------------------------------------
\usepackage{listings}
\usepackage{xcolor}
\definecolor{codebg}{HTML}{F6F8FA}
\definecolor{codekw}{HTML}{0F62FE}
\definecolor{codecmt}{HTML}{6A737D}
\definecolor{codestr}{HTML}{0A7D33}
\definecolor{ipnvino}{HTML}{6B1E3C}   % guinda IPN
\lstset{
  basicstyle=\ttfamily\small,
  backgroundcolor=\color{codebg},
  keywordstyle=\color{codekw}\bfseries,
  commentstyle=\color{codecmt}\itshape,
  stringstyle=\color{codestr},
  numberstyle=\tiny\color{codecmt},
  numbers=left,numbersep=8pt,
  frame=single,rulecolor=\color{codecmt!30},
  breaklines=true,showstringspaces=false,
  tabsize=2,captionpos=b,
  literate={á}{{\'a}}1 {é}{{\'e}}1 {í}{{\'i}}1 {ó}{{\'o}}1 {ú}{{\'u}}1
           {ñ}{{\~n}}1 {Á}{{\'A}}1 {É}{{\'E}}1 {Í}{{\'I}}1 {Ó}{{\'O}}1
           {Ú}{{\'U}}1 {Ñ}{{\~N}}1
}

% ---------- Estilo de títulos (titlesec) ----------------------------
\usepackage{titlesec}
\titleformat{\chapter}[display]
  {\sffamily\huge\bfseries\color{ipnvino}}
  {\sffamily\Large\bfseries\color{ipnvino!70}\MakeUppercase{\chaptertitlename}\ \thechapter}
  {12pt}{\Huge}
\titlespacing*{\chapter}{0pt}{10pt}{30pt}
\titleformat{\section}{\sffamily\Large\bfseries\color{ipnvino}}{\thesection}{1em}{}
\titleformat{\subsection}{\sffamily\large\bfseries}{\thesubsection}{1em}{}
\titleformat{\subsubsection}{\sffamily\normalsize\bfseries}{\thesubsubsection}{1em}{}

% ---------- Encabezados y pies (fancyhdr) ---------------------------
\usepackage{fancyhdr}
\pagestyle{fancy}
\fancyhf{}
\fancyhead[L]{\sffamily\small\nouppercase{\leftmark}}
\fancyhead[R]{\sffamily\small\thepage}
\renewcommand{\headrulewidth}{0.4pt}
\fancypagestyle{plain}{\fancyhf{}\fancyfoot[C]{\sffamily\small\thepage}\renewcommand{\headrulewidth}{0pt}}

% ---------- Bibliografía (biblatex / biber, estilo IEEE) -------------
\usepackage[backend=biber,style=ieee,sorting=none,maxbibnames=6,giveninits=true]{biblatex}
\addbibresource{references.bib}
% Permitir el corte de URLs largas (hashes, rutas) para evitar desbordes
\setcounter{biburlnumpenalty}{9000}
\setcounter{biburlucpenalty}{9000}
\setcounter{biburllcpenalty}{9000}

% ---------- Glosario / Acrónimos ------------------------------------
\usepackage[toc,nonumberlist,nogroupskip,nopostdot]{glossaries}
\makeglossaries

% ---------- Hipervínculos y referencias cruzadas --------------------
\usepackage[hidelinks,bookmarksnumbered,unicode]{hyperref}
\hypersetup{
  pdftitle={Prototipo de un sistema web para la gestión de expedientes digitales y actividad procesal de un despacho legal},
  pdfauthor={Eduardo Alonso Sánchez; Hatziry Vitales Herrera},
  colorlinks=true,linkcolor=ipnvino,citecolor=ipnvino,urlcolor=ipnvino,
}
\usepackage[nameinlink,spanish]{cleveref}
% cleveref en español nombra las secciones «apartado»; usamos «sección»
\crefname{section}{sección}{secciones}
\Crefname{section}{Sección}{Secciones}
\crefname{subsection}{sección}{secciones}
\Crefname{subsection}{Sección}{Secciones}
\crefname{table}{tabla}{tablas}
\Crefname{table}{Tabla}{Tablas}
\crefname{figure}{figura}{figuras}
\Crefname{figure}{Figura}{Figuras}

% ---------- Atajos útiles -------------------------------------------
\newcommand{\figura}[4][\textwidth]{% [width]{file}{caption}{label}
  \begin{figure}[H]\centering
    \includegraphics[width=#1]{#2}
    \caption{#3}\label{#4}
  \end{figure}}
